\section{Newo scrambling and invariance of the lattice criterion}
\label{sec:newo}

The order of a sequence matters when arbitrary prefixes are used. It is
therefore natural to ask whether changing the scalar van der Corput ordering
changes the lattice construction criterion. The following class of digit
permutations gives a precise negative answer. We use Pride's name
\emph{Newo}, Owen spelled backwards, for the reversal of the digit-dependency
direction in nested scrambling~\cite{pride-notes}.

\subsection{Suffix-dependent digit permutations}

Fix $M$ binary digits. At each position $k$, where $0\le k<M$, every
possible suffix fixes a permutation of $\{0,1\}$, denoted by
\[
 \pi_k(\,\cdot\,;a_{k+1},\ldots,a_{M-1}).
\]
For $a=\sum_{k=0}^{M-1}a_k2^{-k-1}$ define
\begin{equation}\label{eq:newo-definition}
 \psi(a)=\sum_{k=0}^{M-1}
 \pi_k(a_k;a_{k+1},\ldots,a_{M-1})2^{-k-1}.
\end{equation}
At $k=M-1$ the suffix is empty. The same collection of permutations is
applied to every input point. Unlike the usual nesting direction, the
permutation of a more significant fractional digit depends on less
significant input digits. One can choose these permutations independently
and uniformly at random, but the invariance result does not require their
randomness.

The map $\psi$ is a bijection of the dyadic grid
$\{0,2^{-M},\ldots,(2^M-1)2^{-M}\}$.
Indeed, given the output, recover the final input digit by inverting its
permutation, then recover preceding digits successively using the suffix
already recovered. Define $t_i=\psi(\phi_2(i))$ for $0\le i<2^M$.
The randomized points considered here are
\begin{equation}\label{eq:newo-points}
 \vx_i'=\fracpart{t_i\vh+\vDelta}.
\end{equation}
Thus the scrambling is scalar and is performed before multiplication by the
fixed integer generating vector. The uniform shift $\vDelta$ is independent
of any random choice of $\psi$.

\subsection{Every prefix has the same expected squared discrepancy}

Here and below the expectation defining $e_n^2$ is over the uniform shift.
For random Newo permutations, the following statement holds conditionally
on every choice of those permutations.

\begin{theorem}[Newo invariance, adapted from Pride's notes]
\label{thm:newo-invariance}
For any fixed integer vector $\vh$, any $1\le n\le2^M$, and any collection
of permutations in~\eqref{eq:newo-definition}, the sequence prefix
\eqref{eq:newo-points} and the unscrambled lattice-sequence prefix have the
same shift-averaged squared periodic-kernel discrepancy. Consequently their
weighted criteria agree for every choice of sample-size weights, including
$\W_N$ for every $N\le2^M$.
\end{theorem}

\begin{proof}
The proof in~\cite{pride-notes} compares differences of the scalar inputs.
For $k\ge0$, let
\[
 S_k=\{(2r+1)2^{-k-1}:0\le r<2^k\}.
\]
Fix $i\ge1$ with binary digits $i_k$, and partition the indices $j<i$ by
the largest bit position $k$ where $i_k\ne j_k$. Such a group is present
exactly when $i_k=1$. Within that group $j_k=0$, all bits above $k$ agree
with those of $i$, and the $k$ lower bits of $j$ are free. Hence the group
has $2^k$ members.

Suffix dependence ensures that the scrambled outputs have identical digits
at positions greater than $k$ and opposite digits at position $k$.
Their difference modulo one therefore lies in $S_k$. Since $\psi$ is
injective, the $2^k$ differences for this group are distinct, so they exhaust
$S_k$. The sets $S_k$ are disjoint. We have proved
\begin{equation}\label{eq:newo-row}
 \{\fracpart{t_i-t_j}:0\le j<i\}
   =\bigcup_{k:i_k=1}S_k.
\end{equation}
The identity permutations give the same right-hand side for the
unscrambled inputs. Thus the earlier-index differences agree as multisets
for every $i$. Their negatives give the reverse ordered differences, and
the diagonal differences are all zero. The complete ordered difference
multisets of every prefix consequently agree.

Multiplication by the integer vector $\vh$ and evaluation of the periodic
kernel preserve equality of the sums over these multisets. Applying
\eqref{eq:shift-esd} proves equality of the shift-averaged squared
discrepancies. Summing with the chosen sample-size weights proves the final
claim.
\end{proof}

The conclusion is an exact identity, not merely equality after averaging
over scrambles. In particular, the objective function supplied to every CBC
step is unchanged. With identical candidate sets and tie-breaking rules,
the construction therefore selects the same generating vector. This
invariance concerns the stated discrepancy criterion; particular integration
errors need not agree, and arbitrary permutations or independently scrambled
coordinates are not covered by the theorem.

\subsection{Gray-code ordering and odd multipliers}

An important example uses
\[
 \pi_k(a_k;a_{k+1},\ldots)=a_k\mathbin{\oplus}a_{k+1},
 \qquad a_M=0,
\]
where $\oplus$ is exclusive-or. If
$g(i)=i\mathbin{\oplus}\lfloor i/2\rfloor$ is the binary reflected Gray
code, then $\psi(\phi_2(i))=\phi_2(g(i))$. The theorem shows why a generating
vector selected using van der Corput ordering has precisely the same
weighted prefix criterion under Gray-code ordering. Consecutive values of
$g(i)$ differ in one bit; arbitrary Newo choices need not have this extra
property.

Multiplication modulo one by an odd integer $h$ also has the form
\eqref{eq:newo-definition}. The $k$th fractional digit of $\fracpart{ha}$ is
\begin{equation}\label{eq:odd-multiplier-newo}
 a_k\mathbin{\oplus}
 \left(\left\lfloor
 h\sum_{\ell=k+1}^{M-1}a_\ell2^{k-\ell}
 \right\rfloor\bmod2\right).
\end{equation}
Indeed, take $\lfloor 2^{k+1}ha\rfloor$ modulo two. Digits above $k$
contribute even integers; the current digit contributes $ha_k$, whose
parity equals $a_k$; the remaining term is the carry from less significant
digits. For each suffix this is either the identity or the swap of the
current digit. Thus ordinary odd multiplication is included within the
Newo class. This does not imply that independent Newo maps in all coordinates
form a rank-one lattice sequence or inherit the same criterion.

\par\smallskip
\noindent\codex{[Author verification: Confirm the finite-precision definition and attribution with Anders.
The odd-multiplier calculation makes explicit the carry argument sketched in
the source notes. The theorem is stated at a fixed maximum precision;
a consistent infinite-precision randomization requires a separate definition.
Newo scrambling does not, in general, preserve the digital-net properties
preserved by ordinary nested scrambling.]}
\par\smallskip
